% titel: Blatt 7, A1: Lewis-Formeln POCl3 und H2SO4 (Oktett + erweitert)
% Teil II: der durchdachte Weg.

\subsection*{Gegeben und gesucht}
$\mathrm{POCl_3}$ und $\mathrm{H_2SO_4}$. F\"ur jede Verbindung zwei Lewis-
Formeln: Oktett erfüllt bzw. erweiterte Valenzschale. Alle Formalladungen
und freien Elektronenpaare angeben.

\subsection*{Ansatz}
Valenzelektronen z\"ahlen, Zentralatom (P bzw. S) w\"ahlen, Gerüst bauen,
Elektronen verteilen, Formalladungen pr\"ufen. Oktett erzwingen bzw. f\"ur
die grosse Form das Zentralatom mit mehr Bindungspaaren versehen.

\subsection*{Rechnung und Zeichnung}
\textbf{a) $\mathrm{POCl_3}$:}\\
Valenzelektronen: P (5) $+$ O (6) $+$ 3$\times$Cl (7) $= 32$ Elektronen
$= 16$ Paare. Zentralatom P, daran O und 3 Cl.

\textbf{Oktett-Form ($\mathrm{POCl_3}$, alle Oktette):}
\nopagebreak
\[
  \mathrm{Cl - \overset{+}{\mathrm{P}} - O^{\,-}}
  \qquad\text{(mit je 3 freien Paaren an jedem Cl und O)}
\]
Aufbau: P bildet nur 3 Einfachbindungen zu Cl und eine zu O (kein P=O),
P beh\"alt 1 freies Elektronenpaar. Dann: O hat 1 Bindung + 3 freie Paare
(Formalladung $-1$); P: $5 - 2\,(\mathrm{freies\,Paar}) - \tfrac{8}{2} = +1$.
Oktett f\"ur alle Atome erf\"ullt.

\textbf{Erweiterte Form ($\mathrm{POCl_3}$):}
\[
  \overset{=}{\mathrm{O}}\mathrm{P(Cl)}_3
  \qquad\text{d.~h.\ P=O-Doppelbindung und 3 P--Cl-Einfachbindungen}
\]
P ist dann von 5 Elektronenpaaren umgeben (1 Doppel + 3 einfach $=$ 5
Bindungspaare $\rightarrow$ 10 Valenzelektronen, erweiterte Valenzschale).
Formalladungen: P $5-0-\tfrac{10}{2}=0$, O $6-4-2=0$, Cl $7-6-1=0$.
Alle Formalladungen sind 0.

\textbf{b) $\mathrm{H_2SO_4}$:}
Valenzelektronen: 2$\times$H (1) $+$ S (6) $+$ 4$\times$O (6) $= 32$
Elektronen $= 16$ Paare. Zentralatom S, 2 OH-Gruppen und 2 O als Enden.

\textbf{Oktett-erfüllende Form ($\mathrm{H_2SO_4}$):}
\[
  \mathrm{H}{-}\mathrm{O}{-}\mathrm{S}(=O)_2{-}\mathrm{O}{-}\mathrm{H}
\]
Damit S nur ein Oktett hat, m\"ussen zwei S=O-Doppelbindungen gezeichnet
werden (S mit 4 Bindungspaaren $=$ 8 Elektronen). Alternativ reicht f\"ur
Oktett nicht aus, alle Bindungen einfach zu lassen w\"urde S ein offenes
Oktett geben. Formalladungen alle 0.

\textbf{Erweiterte Form ($\mathrm{H_2SO_4}$):}
\[
  \mathrm{H}{-}\mathrm{O}{-}\mathrm{S}\overset{\,=O}{\underset{\,=O}{\mathrm{(O)}}}{-}\mathrm{O}{-}\mathrm{H}
  \quad\text{oder S mit Doppelbindung zu 2 O und 2 S--O--H}
\]
S kann auch mehr als 4 Bindungspaare zeigen. Sinnvoll: S bildet
gegen die 2 OH nur Einfachbindungen und gegen 2 O Doppelbindungen; damit
ist die Zahl der Bindungselektronen um S mit 6 Bindungspaaren gr\"osser als
das Oktett $\rightarrow$ erweiterte Valenzschale (S nutzt d-Orbitale).
Formalladungen bleiben 0.

\begin{ergebnis}
\begin{itemize}
  \item $\mathrm{POCl_3}$: Oktett-Form mit P$^{+}$/O$^{-}$ (keine
        Doppelbindung); erweiterte Form mit P=O (alle Formalladungen 0).
  \item $\mathrm{H_2SO_4}$ oktett-erfüllt: O=S=O-Doppelbindungen + 2 OH;
        erweitert: S mit bis zu 6 Bindungspaaren.
\end{itemize}
\end{ergebnis}

\begin{kontrolle}
Beide gezeichneten Strukturen pro Molekül summieren die Valenzelektronen
korrekt (beide 32 bzw. 32 e${}^-$). Die Oktett-erfüllenden Formen tragen
Formalladungen (POCl$_3$: P$^{+}$, O$^{-}$), die erweiterten Formen sind
ladungsneutral und entsprechen der in der Praxis dominanten Lewis-Struktur.
\end{kontrolle}

\begin{merksatz}
Atome der 3. Periode (P, S) ``brechen'' die Oktettregel und k\"onnen mehr
als 4 Bindungspaare aufnehmen. Zeichnet man eine Oktett-erf\"ullende Form,
muss man Formalladungen einzeichnen; die erweiterte Form mit
Doppelbindungen (oder 6 Bindungspaaren) hat meist alle Formalladungen 0.
\end{merksatz}
